\addcontentsline{toc}{chapter}{Abstract}

\begin{center}
\huge{{\bf Abstract}}
\vspace{2cm}
\end{center}

MicroTEDAClus is an evolving clustering algorithm for data streams that summarizes the stream into microclusters by the eccentricity and the typicality of the TEDA framework, and whose authors report that it performs on par with established stream clustering algorithms and adapts better than them to changes in the data. This undergraduate thesis, developed at the Federal University of Minas Gerais in the Bachelor's Degree in Control and Automation Engineering, under the supervision of Prof. Janier Arias Garcia, reassesses the algorithm and proposes an improved model. The reassessment finds issues in the formulation of the algorithm, differences between the formulation and the public implementation that produced its published results, and flaws in the methodology of its experiments, among them preprocessing steps that leak information about the data and an evaluation that predicts each batch of samples after learning it. Under a methodology without these flaws, in which every model is executed in the same language on the unprocessed data, evaluated prequentially, tuned by the same automatic search and compared over many runs by nonparametric statistical tests, MicroTEDAClus is frequently outperformed by the benchmark models, and its claimed adaptability is not supported in full.

The improved model, MicroTEDAClus+, results from controlled experiments in which each correction or modification is accepted by a predefined rule, on runs separate from those of its evaluation. On the synthetic datasets of the original work it improves on MicroTEDAClus in clustering quality, execution time and memory, it is among the best benchmark models in the datasets with concept drift and the fastest of them, and in an experiment modernized in its environment, its benchmark models and its selection of hyperparameters it remains among the best models in most datasets. On eight real datasets, however, it is among the worst models. A diagnosis and a second effort of controlled experiments on those datasets produce the final version of MicroTEDAClus+, which lowers the search range of its variance limit, drops its noise label and searches the factor that links its microclusters. It is no longer among the worst models in most real datasets and is the best in one of them, but it remains behind the best benchmark models on the real streams with concept drift.

\clearpage
\thispagestyle{empty}
\cleardoublepage
